\section{Results}
\label{sec:results}

\subsection{Effects during the contract period}
Table~\ref{tab:main} reports ITT estimates of equation~\eqref{eq:main}. The estimates use all 2,416 randomized employees and their weekly attendance in the eight intervention weeks (columns 1 and 2) and in the twelve weeks that followed (columns 3 and 4).

We focus on column (2), which reports the primary outcome, the number of visits per week. Offering a contract without reminders raised attendance by 0.31 visits per week (95 percent confidence interval 0.17 to 0.45), from 0.84 to 1.15 visits per week, an increase of 37 percent. The point estimate for reminders alone is 0.064 visits per week (standard error 0.052). This estimate is small and not statistically significant at conventional levels. The interaction term is positive: reminders raise the effect of the contract offer by 0.152 visits per week. Employees offered both a contract and reminders made 0.51 more visits per week than the control group, an increase of 57 percent.

Because the regression includes the interaction, the coefficient on the contract offer measures its effect in the absence of reminders, and the coefficient on reminders measures their effect in the absence of an offer. The effect of the offer when reminders are also sent is the sum of the first and third coefficients, 0.462 visits per week.

Column (1) shows the same pattern for the share of weeks with at least one visit. The contract offer raised this share by 6.5 percentage points, or 17 percent of the control mean of 38.0 percent. The estimates for reminders alone and for the interaction are small and imprecise for this outcome.

\begin{table}[t]
\centering
\small
\begin{threeparttable}
\caption{Effects of contract offers and reminders on gym attendance}
\label{tab:main}
\begin{tabular}{@{}p{5.7cm}cccc@{}}
\toprule
 & (1) & (2) & (3) & (4) \\
 & \multicolumn{2}{c}{Weeks 1--8} & \multicolumn{2}{c}{Weeks 9--20} \\
 & Any visit & Visits per week & Any visit & Visits per week \\
\midrule
Commitment offer & $0.065^{***}$ & $0.310^{***}$ & $0.019$ & $0.112^{**}$ \\
 & (0.017) & (0.071) & (0.013) & (0.059) \\
Reminders & $0.012$ & $0.046$ & $-0.003$ & $0.009$ \\
 & (0.015) & (0.052) & (0.012) & (0.044) \\
Commitment offer × Reminders & $0.024$ & $0.152^{**}$ & $0.011$ & $0.038$ \\
 & (0.024) & (0.076) & (0.019) & (0.064) \\
\midrule
Control mean & 0.380 & 0.84 & 0.335 & 0.71 \\
Observations (participant-weeks) & 19,296 & 19,296 & 28,944 & 28,944 \\
Participants & 2,412 & 2,412 & 2,412 & 2,412 \\
\bottomrule
\end{tabular}
\begin{tablenotes}
\small
\item[] Notes: Intention-to-treat estimates of equation~\eqref{eq:main} from linear regressions of weekly outcomes on the arm indicators, week fixed effects and randomization-stratum fixed effects. Any visit is an indicator for at least one visit in the week, so its mean is the share of weeks with a visit. Weeks 1--8 are the contract period and weeks 9--20 the twelve weeks after it. Control mean is the mean of the outcome in the control arm over the same weeks. Standard errors clustered by participant in parentheses. * $p<0.10$, ** $p<0.05$, *** $p<0.01$.
\end{tablenotes}
\end{threeparttable}
\end{table}


\subsection{After the contracts expired}
Columns (3) and (4) cover weeks 9 to 20. The effect of the contract offer on visits falls to 0.112 per week, about a third of its size during the contract period, and remains significant at the 5 percent level. The share of weeks with a visit is 1.9 percentage points higher in the commitment arm, with a 95 percent interval of 0.4 to 3.4 points. The reminder estimates are close to zero after week 8, and so are the interaction terms.

\subsection{Signers and decliners}
How much of the effect comes from the contracts themselves? During the eight contract weeks, signers made 1.94 visits per week and decliners 0.89. Signing a contract therefore more than doubled gym attendance. Signers had also been more frequent gym users before the study, with 1.18 visits per week at baseline against 0.70 for decliners. Random assignment offers a second way to scale the effect: dividing the effect of the offer in the arm without reminders (0.310) by the take-up rate in that arm (32.8 percent) gives about 0.95 visits per week for those who signed.

\subsection{Who signs a contract}
Take-up was higher among employees who said at baseline that they exercised less often than their own target: 37.1 percent against 27.3 percent. This pattern is consistent with \citet{royer2015}, who find that the demand for commitment comes mainly from employees who fall short of their own exercise targets. Women were more likely than men to sign (36.3 against 30.4 percent). These take-up correlates were not pre-specified and are descriptive.

\subsection{Heterogeneity}
Table~\ref{tab:het} reports the effect of the contract offer in five pairs of subgroups, pooling the arms with and without reminders. The pooled effect is 0.386 visits per week, the average of the effects of the offer with and without reminders in Table~\ref{tab:main}. The commitment offer works mainly for employees who rarely used the gym before the study: among employees with fewer than one weekly visit at baseline, it raised visits by 0.508 per week, against 0.187 among more frequent users (p = 0.040 for the difference). The other four splits, by gender, age, position and travel time to the gym, show no significant differences.

\begin{table}[t]
\centering
\small
\begin{threeparttable}
\caption{Effect of the contract offer on visits per week by subgroup, weeks 1--8}
\label{tab:het}
\begin{tabular}{@{}p{6.4cm}ccc@{}}
\toprule
 & (1) & (2) & (3) \\
Split (group A versus group B) & Group A & Group B & p-value, A = B \\
\midrule
Women versus men & $0.431^{***}$ & $0.316^{***}$ & 0.36 \\
 & (0.078) & (0.098) & \\
Under 40 versus 40 and over & $0.352^{***}$ & $0.415^{***}$ & 0.61 \\
 & (0.092) & (0.083) & \\
Faculty versus staff and other & $0.288^{**}$ & $0.425^{***}$ & 0.33 \\
 & (0.121) & (0.071) & \\
Fewer than one baseline visit per week versus one or more & $0.508^{***}$ & $0.187$ & 0.040 \\
 & (0.088) & (0.129) & \\
Under 10 minutes to the gym versus 10 minutes or more & $0.398^{***}$ & $0.371^{***}$ & 0.83 \\
 & (0.081) & (0.094) & \\
\midrule
All participants & $0.386^{***}$ & & \\
 & (0.061) & & \\
\bottomrule
\end{tabular}
\begin{tablenotes}
\small
\item[] Notes: Each row reports the effect of the contract offer, pooling the arms with and without reminders, estimated separately in the two groups of a split, with the specification of Table~\ref{tab:main}. Column (3) reports the p-value of a test that the two effects are equal; p-values are not adjusted for multiple comparisons. Standard errors clustered by participant in parentheses. * $p<0.10$, ** $p<0.05$, *** $p<0.01$.
\end{tablenotes}
\end{threeparttable}
\end{table}

